\subsection{主要工作总结}

本文围绕 LeRobot 在线控制中的``感知--推理--执行''闭环，研究了树莓派 5
这类无 GPU ARM 边缘平台上机器人基础软件栈的系统级性能问题。与只统计端到端
FPS 或只从模型、硬件通信单点优化入手的做法不同，本文将问题拆解到 AI 推理框架层、Python 运行时层、Linux
内核层与硬件 I/O 层，建立了面向机器人在线控制负载的四层性能剖析方法。
该方法能够把上层策略调用、运行时线程行为、系统调用阻塞和底层设备通信放到同一条
时间线上观察，从而为后续优化判断提供依据。

在实验环境方面，本文以 SO-101 六自由度机械臂、双摄像头输入和 ACT 策略模型为对象，
构建了可复现实验平台，并围绕 \texttt{record.py} 中的核心控制循环记录各阶段耗时。
在测量工具层面，本文比较了 Python 日志、\texttt{torch.profiler}、
\texttt{strace} 与 \texttt{ftrace} 的适用边界。工具链校准结果表明，
\texttt{strace} 会显著放大上下文切换并改变阶段时间结构，因此不适合承担定量时序分析；
\texttt{ftrace} 延迟导出方案对主循环扰动相对可控，更适合作为内核态观测工具。
在此基础上，本文进一步通过自定义内核插桩拆分了 \texttt{pselect6} 和
\texttt{futex} 路径中的阻塞时间与 CPU 处理时间，使用户态日志和内核态事件能够互相印证。

性能剖析结果显示，当前系统的主要瓶颈并不在单个系统调用或舵机写入函数本身，
而在 ARM CPU 上执行 ACT 模型前向推理的绝对耗时。感知阶段中，两路摄像头已经通过
OpenCV 后台线程完成异步预取，主线程读取缓存帧的开销较小；执行阶段中，
舵机 \texttt{sync\_read} 与 \texttt{sync\_write} 的毫秒级耗时主要来自
1 Mbps 半双工总线和 SCServo 协议约束，而不是 Linux 系统调用路径。
相比之下，单次完整 ACT 推理耗时约为秒级，其中 ResNet18 视觉特征提取和
Transformer Encoder 占据主要比例，是限制在线控制频率的核心因素。
这一结论说明，在当前硬件条件下，优化方向应优先指向推理运行时和模型算子执行效率，
而不是盲目改写已受物理总线约束的硬件 I/O 路径。

在优化探索方面，本文首先分析了感知、执行与推理三类路径的异步化可行性。
由于舵机读写共享同一条半双工总线，简单后台线程化会退化为``持锁串行 + 线程切换''，
难以形成真正并行；相机路径已经具备后台预取机制，进一步收益有限。
因此，本文将优化重点放在推理异步化上，推导了异步流水线稳定无 stall 的近似条件
\(T_{\text{infer}}F \leq \min(hN,\;N/2)\)，并进一步建立了
FPS 随 hold 在临界点前上升、临界点后饱和到目标帧率的分段模型。
实验结果表明，在 \texttt{chunk\_size=100}、目标 30 FPS 的设置下，
推荐配置 hold=0.30 可取得 27.42 FPS 的有效控制频率，相对同步 baseline
提升 70.0\%，且稳态阶段几乎无 stall，说明异步调度能够显著回收动作执行期间的空窗。
episode 平均 FPS 尚未达到 30 的主要原因是启动空队列等待和树莓派 5 单机部署下的 CPU 资源争抢，
后续仍需从机器人基础软件栈中的推理运行时、操作系统调度和跨层诊断工具继续优化。

\subsection{研究局限性}

本文的实验对象固定在树莓派 5、SO-101 机械臂和 ACT 策略模型上，
因此结论首先服务于 ARM CPU 边缘机器人这一类受限平台。
不同机械臂的通信协议、相机驱动、策略模型结构和任务动作尺度可能改变各阶段占比，
尤其是在模型更小、舵机总线更慢或相机分辨率更高的配置中，
瓶颈可能从推理运行时转移到硬件 I/O 或图像预处理路径。
因此，本文得到的具体数值不应直接外推到所有机器人平台，
更适合作为一套跨层性能分析方法和基础软件优化思路的案例验证。

其次，本文已经对异步推理进行了实测验证，但对推理运行时本体的优化仍然有限。
ONNX Runtime、INT8 量化、算子融合、ARM NEON 后端优化以及 PyTorch 线程池参数等方向，
在本文中主要作为瓶颈定位后的推论和未来工作提出，尚未形成完整的对比实验。
因此，本文能够说明当前系统为什么慢、异步调度能恢复多少吞吐，
但还不能给出降低单次 ACT 推理耗时的最终工程方案。

最后，本文的动作质量评价仍以操作者现场观察为主，尚未建立独立于主观判断的任务质量指标。
异步推理中的 stall 比例、过期帧比例和有效 FPS 可以反映控制流水线的运行状态，
但它们并不能完全替代抓取成功率、末端轨迹误差、动作平滑度和安全裕量等机器人任务指标。
后续若要将本文方法推广为更通用的基础软件评测框架，
仍需把系统级性能指标与任务级行为质量进一步关联起来。

\subsection{未来工作}

未来工作应继续围绕机器人基础软件栈展开，而不是停留在简单更换任务或平台的层面。
首先，在推理运行时方面，需要针对 ARM CPU 上 ACT 推理耗时过高的问题开展系统优化实验。
可优先比较 PyTorch eager 模式、ONNX Runtime、INT8 量化和算子融合方案在树莓派 5
上的端到端收益，并进一步分析 ResNet18 卷积、Transformer 线性层和 GEMM 算子
在 NEON 或 KleidiAI 等 ARM 后端上的执行效率。同时，还需要评估 PyTorch
线程池规模、OpenMP 调度策略和模型 batch/shape 固定化对延迟抖动的影响，
使推理运行时优化能够服务于稳定控制频率，而不仅是单次 benchmark 的峰值速度。

其次，在操作系统与调度层面，可以继续研究 Linux 对机器人在线控制负载的支持能力。
当前异步推理实验中，\texttt{policy\_server}、\texttt{robot\_client}、
相机线程和串口通信都竞争同一组 CPU 核心，因而推理耗时和控制循环 jitter
会受到调度策略影响。后续可围绕 CPU 频率锁定、线程亲和性、实时优先级、
\texttt{cgroup} 资源隔离和 PREEMPT\_RT 内核展开对比，量化这些机制对
\texttt{futex} 等待、上下文切换、队列 stall 和有效 FPS 的影响。
这类实验有助于回答机器人基础软件中一个更本质的问题：
在有限 CPU 资源上，推理任务、控制任务和 I/O 任务应如何被操作系统共同调度。

再次，在硬件 I/O 基础库方面，仍需区分``可以被软件隐藏的等待''和
``由协议物理约束决定的等待''。摄像头路径可继续从 V4L2 缓冲队列、
OpenCV 后台线程、MJPEG 解码和图像张量转换等层面分析端到端时延，
明确后台预取在不同帧率与分辨率下的边界。舵机路径则应围绕 SCServo 半双工总线、
串口波特率、同步读写协议和潜在的流水线读机制继续建模，
判断哪些优化只是移动阻塞位置，哪些优化能够真正减少主控制路径耗时。
这种面向 I/O 基础库的分析对机器人系统尤其重要，因为硬件通信往往同时涉及性能、
安全和动作语义，不能只按普通服务器程序中的异步 I/O 经验处理。

最后，本文使用的跨层 profiling 方法仍可进一步工具化。
后续可以将 Python 日志、\texttt{torch.profiler}、\texttt{ftrace}、
内核自定义插桩和实验 CSV 自动对齐为统一时间线，
形成面向机器人基础软件的可复用性能诊断工具。
该工具应能够自动标注每一帧中的感知、推理、执行和等待区间，
关联用户态函数、内核阻塞点和硬件设备事件，
并输出适合优化决策的指标，例如推理临界时间、stall 触发原因、
I/O 等待占比和调度抖动来源。通过这种方式，本文的工作可以从一次针对 LeRobot
的系统剖析，进一步发展为面向边缘机器人基础软件栈的通用评测与优化方法。
